% Zwischenpräsentation EAD-Projekt: TCGA PanCanAtlas
% Übersetzen mit:  latexmk -lualatex zwischenpraesentation.tex   (im Ordner praesentation/)
% Abbildungen erzeugen (aus dem Projektordner):  .venv/bin/python praesentation/abbildungen.py
\documentclass[aspectratio=169,11pt]{beamer}

\usepackage{fontspec}
\usefonttheme{professionalfonts}
\setsansfont{Fira Sans}[UprightFont=*-Regular, BoldFont=*-SemiBold, ItalicFont=*-Italic, BoldItalicFont=*-SemiBoldItalic]
\usepackage[ngerman]{babel}
\usepackage{amsmath,amssymb}
\usepackage{booktabs,tabularx,array}
\usepackage{tikz}
\usetikzlibrary{positioning,calc,arrows.meta,decorations.pathreplacing,shapes.misc}
\usepackage[most]{tcolorbox}

% ------------------------------------------------------------------ Farben (aus der Vorlage slides-handout.pdf)
\definecolor{titelbalken}{HTML}{344956}
\definecolor{text}{HTML}{17212B}
\definecolor{grau}{HTML}{5C6872}
\definecolor{linie}{HTML}{D3D8DC}
\definecolor{petrol}{HTML}{006460}
\definecolor{petrolhell}{HTML}{F2F8F8}
\definecolor{blau}{HTML}{005F9E}
\definecolor{blaurand}{HTML}{669FC5}
\definecolor{blauhell}{HTML}{F2F7FA}
\definecolor{orange}{HTML}{E9A03B}
\definecolor{dunkel}{HTML}{1C2A36}

% ------------------------------------------------------------------ Beamer-Vorlagen
\setbeamertemplate{navigation symbols}{}
\setbeamercolor{normal text}{fg=text}
\setbeamercolor{alerted text}{fg=blau}
\setbeamercolor{structure}{fg=titelbalken}
\setbeamerfont{frametitle}{size=\large}
\setbeamertemplate{frametitle}{%
  \nointerlineskip
  \begin{beamercolorbox}[wd=\paperwidth,ht=2.6ex,dp=1.1ex,leftskip=0.5cm]{frametitle}
    \usebeamerfont{frametitle}\insertframetitle
  \end{beamercolorbox}}
\setbeamercolor{frametitle}{bg=titelbalken,fg=white}
\setbeamertemplate{footline}{%
  \hfill\usebeamerfont{footline}\color{grau}\scriptsize\insertframenumber\hspace*{0.45cm}\vspace*{0.3cm}}
\setbeamertemplate{itemize item}{\raisebox{0.15ex}{\color{titelbalken}\textbf{--}}}
\setbeamertemplate{itemize subitem}{\color{grau}\textbf{--}}
\setbeamertemplate{enumerate item}{\color{titelbalken}\textbf{\insertenumlabel.}}
\setlength{\leftmargini}{1.2em}
\setbeamersize{text margin left=0.9cm, text margin right=0.9cm}

% ------------------------------------------------------------------ Kästen wie in der Vorlage
% Kasten mit hellblauem Grund und blauem Rand, optional mit Überschrift
\newtcolorbox{kasten}[1][]{enhanced, colback=blauhell, colframe=blaurand, boxrule=0.5pt, arc=2pt,
  left=5pt, right=5pt, top=3pt, bottom=3pt, fontupper=\color{text}\raggedright,
  #1}
\newcommand{\kastentitel}[1]{{\color{blau}\textbf{#1}}\par\smallskip}
% Beispiel-Kasten: Petrol-Titelzeile, heller Grund
\newtcolorbox{beispiel}[2][]{enhanced, colback=petrolhell, colframe=petrolhell, coltitle=white, colbacktitle=petrol,
  boxrule=0pt, arc=0pt, left=5pt, right=5pt, top=3pt, bottom=3pt, toptitle=1.5pt, bottomtitle=1.5pt,
  title={\textbf{Beispiel}\enspace #2}, #1}

\newcommand{\gen}[1]{\texttt{#1}}
\newcommand{\quelle}[1]{\par\vspace{2pt}{\color{grau}\tiny #1}}
\renewcommand{\arraystretch}{1.15}
\newcolumntype{Y}{>{\raggedright\arraybackslash}X}

% ------------------------------------------------------------------ Titelseite und Abschnittsseiten
\newcommand{\titelgrafik}{%
  % kleine Expressionsmatrix: Zeilen = Proben (links nach Tumortyp eingefärbt), Spalten = Gene
  \begin{tikzpicture}[x=0.27cm,y=0.27cm]
    \pgfmathsetseed{7}
    \foreach \r in {0,...,13} {
      \pgfmathtruncatemacro{\gruppe}{\r/5}
      \ifcase\gruppe \def\gf{blau}\or \def\gf{petrol}\else \def\gf{orange}\fi
      \fill[\gf] (-1.2,-\r) rectangle ++(0.75,-0.9);
      \foreach \c in {0,...,9} {
        \pgfmathsetmacro{\basis}{(\gruppe==0 && \c<3) || (\gruppe==1 && \c>3 && \c<7) || (\gruppe==2 && \c>6) ? 70 : 12}
        \pgfmathsetmacro{\w}{max(4,min(95,\basis+22*rnd-11))}
        \fill[blau!\w!white] (\c,-\r) rectangle ++(0.9,-0.9);
      }
    }
    \node[anchor=south,font=\scriptsize,text=grau] at (4.95,0.15) {Gene};
    \node[anchor=south,font=\scriptsize,text=grau,rotate=90] at (-1.5,-6.95) {Proben};
  \end{tikzpicture}}

\setbeamertemplate{title page}{%
  \begin{minipage}[c]{0.66\textwidth}\raggedright
    {\usebeamerfont{title}\Large\bfseries\color{text}\inserttitle\par}
    \vspace{0.45cm}
    {\color{grau}\insertsubtitle\par}
    \vspace{0.35cm}
    {\color{linie}\rule{\linewidth}{0.5pt}}\par
    \vspace{0.4cm}
    {\small\insertauthor\par\vspace{0.15cm}\color{grau}\insertinstitute\par\vspace{0.1cm}\insertdate\par}
  \end{minipage}\hfill
  \begin{minipage}[c]{0.3\textwidth}\centering\titelgrafik\end{minipage}}

\newcommand{\abschnittsfolie}[1]{%
  \begin{frame}[plain,noframenumbering]
    \begin{tikzpicture}[remember picture,overlay]
      \fill[dunkel] ($(current page.north east)+(-0.42\paperwidth,0)$) rectangle (current page.south east);
      \begin{scope}[shift={($(current page.south east)+(-0.42\paperwidth,0)$)}]
        \pgfmathsetseed{3}
        \foreach \i in {1,...,9} \draw[white,opacity=0.06] (0,\i*0.95) -- (0.42\paperwidth,\i*0.95);
        \foreach \i in {1,...,7} \draw[white,opacity=0.06] (\i*0.95,0) -- (\i*0.95,\paperheight);
        \foreach \i in {1,...,14} {
          \pgfmathsetmacro{\px}{0.4+5.9*rnd}\pgfmathsetmacro{\py}{0.4+8.2*rnd}
          \pgfmathtruncatemacro{\f}{3*rnd}
          \ifcase\f \def\pf{orange}\or \def\pf{blaurand}\else \def\pf{petrolhell}\fi
          \fill[\pf,opacity=0.85] (\px,\py) circle (1.3pt);
          \fill[\pf,opacity=0.15] (\px,\py) circle (4pt);
        }
      \end{scope}
    \end{tikzpicture}
    \begin{minipage}{0.5\textwidth}
      {\LARGE\bfseries\color{text}#1\par}
      \vspace{0.5cm}{\color{linie}\rule{\linewidth}{0.5pt}}
    \end{minipage}
  \end{frame}}


% ------------------------------------------------------------------ Entscheidungen und Abbildungen
% Entscheidungs-Kasten: Petrol-Titelzeile „Entscheidung“
\newtcolorbox{entscheidung}[1][]{enhanced, colback=petrolhell, colframe=petrol, boxrule=0.6pt, arc=2pt, coltitle=white, colbacktitle=petrol,
  left=6pt, right=6pt, top=4pt, bottom=4pt, toptitle=2pt, bottomtitle=2pt, fonttitle=\bfseries,
  title={Entscheidung}, fontupper=\color{text}\raggedright, #1}
\newcommand{\bild}[2][\textwidth]{{\centering\includegraphics[width=#1]{abbildungen/#2}\par}}
\newcommand{\gross}[1]{{\Large\bfseries\color{petrol}#1}}
% Kennzahl: große Zahl mit kurzer Beschriftung
\newcommand{\kennzahl}[2]{\begin{minipage}[t]{0.3\textwidth}\centering{\LARGE\bfseries\color{blau}#1}\par\vspace{2pt}{\scriptsize\color{grau}#2}\end{minipage}}

% ------------------------------------------------------------------ Angaben
\title{Wie ähneln sich Tumortypen\\in ihrer Genexpression?}
\subtitle{Zwischenpräsentation: Datensatz, Aufbereitung, Hauptkomponentenanalyse und UMAP}
\author{Marinus Graf}
\institute{Effiziente Algorithmen in Data Science und KI}
\date{Wintersemester 2026/27}

\begin{document}

\begin{frame}[plain,noframenumbering]
  \titlepage
\end{frame}

% ==================================================================
\abschnittsfolie{1.\enspace Der Datensatz}

\begin{frame}{Herkunft der Daten}
  \centering
  \begin{tikzpicture}[>={Stealth[length=5pt]},
      stufe/.style={draw=blaurand,fill=blauhell,rounded corners=3pt,minimum width=3.9cm,minimum height=2.6cm,
                    align=center,inner sep=4pt},
      institut/.style={draw=petrol!60,fill=petrolhell,rounded corners=3pt,minimum width=3.9cm,minimum height=1.5cm,
                    align=center,inner sep=4pt},
      pfeil/.style={->,grau,line width=0.9pt}]
    \node[institut] (nci) at (0,0.95) {{\bfseries\footnotesize National Cancer Institute}\\[1pt]
      {\scriptsize\color{grau}Krebsforschung (USA)}};
    \node[institut] (nhgri) at (0,-0.95) {{\bfseries\footnotesize National Human Genome}\\{\bfseries\footnotesize Research Institute}\\[1pt]
      {\scriptsize\color{grau}Genomforschung (USA)}};
    \node[stufe] (tcga) at (5.6,0) {{\color{blau}\bfseries TCGA}\\[1pt]
      {\scriptsize The Cancer Genome Atlas}\\[3pt]
      {\scriptsize\color{grau}2006--2018}\\[2pt]
      {\footnotesize 11.000 Tumoren aus}\\{\footnotesize 33 Krebsarten vermessen}};
    \node[stufe] (pan) at (10.1,0) {{\color{blau}\bfseries PanCanAtlas}\\[1pt]
      {\scriptsize Pan-Cancer Atlas}\\[3pt]
      {\scriptsize\color{grau}2018}\\[2pt]
      {\footnotesize alle Krebsarten gemeinsam,}\\{\footnotesize einheitlich aufbereitet}};
    \coordinate (zusammen) at (2.35,0);
    \draw[grau,line width=0.9pt] (nci.east) -| (zusammen);
    \draw[grau,line width=0.9pt] (nhgri.east) -| (zusammen);
    \draw[pfeil] (zusammen) -- (tcga.west);
    \node[font=\scriptsize\bfseries,text=grau,above] at (3.0,0) {starten};
    \draw[pfeil] (tcga) -- (pan);
  \end{tikzpicture}
\end{frame}

\begin{frame}{Genutzte Daten aus dem PanCanAtlas}
  \centering
  \begin{tikzpicture}[>={Stealth[length=5pt]},
      quelle/.style={draw=blaurand,fill=blauhell,rounded corners=3pt,minimum width=4.4cm,minimum height=1.4cm,
                    align=center,inner sep=4pt},
      datei/.style={draw=petrol!60,fill=petrolhell,rounded corners=3pt,minimum width=4.2cm,minimum height=2.3cm,
                    align=center,inner sep=4pt},
      pfeil/.style={->,grau,line width=0.9pt}]
    \node[quelle] (pan) at (4.6,0) {{\color{blau}\bfseries PanCanAtlas}\\[1pt]
      {\scriptsize\color{grau}einheitlich aufbereitete Dateien aller 33 Krebsarten}};
    \node[datei] (f1) at (0,-3) {{\bfseries\footnotesize Genexpression}\\[2pt]
      {\scriptsize Sequenzierung der RNA}\\[4pt]
      {\footnotesize 20.531 Gene $\times$}\\{\footnotesize 11.069 Proben}};
    \node[datei] (f2) at (4.6,-3) {{\bfseries\footnotesize Klinische Tabelle}\\[2pt]
      {\scriptsize eine Zeile pro Patient}\\[4pt]
      {\footnotesize Tumortyp, Geschlecht,}\\{\footnotesize Alter, Histologie}};
    \node[datei] (f3) at (9.2,-3) {{\bfseries\footnotesize Qualitätsangaben}\\[2pt]
      {\scriptsize eine Zeile pro Probe}\\[4pt]
      {\footnotesize Sequenziergerät,}\\{\footnotesize Markierung „nicht verwenden“}};
    \draw[grau,line width=0.9pt] (pan.south) -- (4.6,-1.4);
    \draw[grau,line width=0.9pt] (0,-1.4) -- (9.2,-1.4);
    \foreach \f/\x in {f1/0,f2/4.6,f3/9.2} \draw[pfeil] (\x,-1.4) -- (\f.north);
  \end{tikzpicture}
\end{frame}

% ------------------------------------------------------------------
\begin{frame}{Aufbau: eine große Matrix}
  \begin{columns}[c]
    \begin{column}{0.5\textwidth}
      \centering
      \begin{tikzpicture}[>={Stealth[length=4pt]}]
        \fill[blauhell] (0,0) rectangle (4.4,-2.6);
        \draw[blaurand] (0,0) rectangle (4.4,-2.6);
        \foreach \y in {-0.65,-1.3,-1.95} \draw[blaurand!40] (0,\y) -- (4.4,\y);
        \foreach \x in {1.1,2.2,3.3} \draw[blaurand!40] (\x,0) -- (\x,-2.6);
        \fill[blau] (2.2,-1.3) rectangle (3.3,-1.95);
        \node[font=\scriptsize,text=white] at (2.75,-1.625) {$x_{ij}$};
        \draw[decorate,decoration={brace,amplitude=4pt},grau] (0,0.12) -- (4.4,0.12)
          node[midway,above=4pt,font=\footnotesize,text=text] {20.531 Gene (Spalten)};
        \draw[decorate,decoration={brace,amplitude=4pt,mirror},grau] (-0.12,0) -- (-0.12,-2.6)
          node[midway,left=5pt,font=\footnotesize,text=text,align=right] {11.069\\Proben\\(Zeilen)};
        \fill[petrolhell] (0,-3.0) rectangle (4.4,-3.6);
        \draw[petrol!50] (0,-3.0) rectangle (4.4,-3.6);
        \node[font=\scriptsize,text=petrol] at (2.2,-3.3) {Tumortyp, Geschlecht, Alter, \dots};
        \draw[<->,grau] (2.2,-2.65) -- (2.2,-2.95) node[midway,right,font=\scriptsize,text=grau] {verknüpft über den Probennamen};
      \end{tikzpicture}
    \end{column}
    \begin{column}{0.46\textwidth}
      \begin{itemize}
        \item \textbf{Zeile:} eine Gewebeprobe
        \item \textbf{Spalte:} ein Gen
        \item \textbf{Eintrag $x_{ij}$:} wie aktiv Gen $j$ in Probe $i$ ist
        \item Tumortyp und klinische Angaben stehen in den beiden anderen Dateien
      \end{itemize}
      \vspace{0.2cm}
      \begin{kasten}
        Die Matrix ist \textbf{breit}: fast doppelt so viele Gene wie Proben.
      \end{kasten}
    \end{column}
  \end{columns}
\end{frame}

% ------------------------------------------------------------------
\begin{frame}{Beispiel: zwei Zeilen der Matrix}
  \centering
  \vspace*{0.3cm}
  \begin{tikzpicture}[
      zelle/.style={draw=blaurand,fill=blauhell,minimum width=1.5cm,minimum height=0.85cm,align=center,
                    font=\normalsize\bfseries,text=blau},
      hervor/.style={draw=petrol!70,fill=petrolhell,text=petrol},
      kopf/.style={minimum width=1.5cm,font=\footnotesize\bfseries,anchor=south},
      zeile/.style={anchor=east,font=\footnotesize},
      erkl/.style={align=center,font=\scriptsize,text=grau,anchor=north}]
    \def\w{1.5}
    \foreach \i/\gen/\text in {
        0/GAPDH/{überall\\aktiv},
        1/XIST/{nur bei\\Frauen},
        2/RPS4Y1/{nur bei\\Männern},
        3/SFTPC/{Lungen-\\bläschen},
        4/NKX2-1/{Lungen-\\gewebe},
        5/ALB/{Albumin,\\Leber}} {
      \node[kopf] at (\w*\i,0.425) {\gen};
      \node[erkl] at (\w*\i,-1.4) {\text};
    }
    \node[kopf,minimum width=0.9cm,text=grau] at (8.7,0.425) {\dots};
    \node[erkl] at (8.7,-1.4) {20.525\\weitere};
    \node[zeile] at (-0.8,0) {TCGA-05-4402-\dots};
    \foreach \i/\wert/\stil in {0/85.894/,1/3.458/hervor,2/{1,2}/hervor,3/1.610/,4/2.779/,5/{25,6}/}
      \node[zelle,\stil] at (\w*\i,0) {\wert};
    \node[zeile] at (-0.8,-0.85) {TCGA-05-4249-\dots};
    \foreach \i/\wert/\stil in {0/22.878/,1/{1,9}/hervor,2/3.411/hervor,3/2.417/,4/4.180/,5/{1,6}/}
      \node[zelle,\stil] at (\w*\i,-0.85) {\wert};
    \foreach \y in {0,-0.85}
      \node[zelle,minimum width=0.9cm,draw=blaurand!60,dashed,fill=white,text=grau,font=\footnotesize] at (8.7,\y) {\dots};
    \draw[decorate,decoration={brace,amplitude=4pt},grau] (-0.75,1.0) -- (9.15,1.0)
      node[midway,above=5pt,font=\footnotesize\bfseries,text=text] {Gene (Spalten)};
    \draw[decorate,decoration={brace,amplitude=4pt,mirror},grau] (-3.3,0.425) -- (-3.3,-1.275)
      node[midway,left=6pt,font=\footnotesize\bfseries,text=text,align=right] {Proben\\(Zeilen)};
  \end{tikzpicture}

  \vspace{0.15cm}
  {\scriptsize\color{grau}beide: Lungenkrebs (Adenokarzinom der Lunge, LUAD)}
\end{frame}

% ------------------------------------------------------------------
\begin{frame}{Der Probenname (Barcode)}
  \centering
  \begin{tikzpicture}[
      teil/.style={draw=blaurand,fill=blauhell,minimum width=1.38cm,minimum height=0.7cm,
                   font=\normalsize\ttfamily\bfseries,text=blau},
      bed/.style={align=center,font=\scriptsize,anchor=north,text width=2.2cm},
      klammer/.style={decorate,decoration={brace,amplitude=4pt},grau},
      gruppe/.style={midway,above=5pt,font=\footnotesize\bfseries,text=text}]
    \def\d{1.5}
    \foreach \i/\a in {0/TCGA,1/05,2/4402,3/01,4/A,5/01,6/R,7/1206,8/07}
      \node[teil] (o\i) at (\d*\i,0) {\a};
    \foreach \i/\text in {
        0/{Projekt},
        1/{Klinik},
        2/{Patient},
        3/{Probentyp},
        4/{Röhrchen},
        5/{Teilstück},
        6/{Molekül},
        7/{Platte},
        8/{Zentrum}}
      \node[bed] at ([yshift=-0.15cm]o\i.south) {\text};
    \draw[klammer] ([yshift=0.12cm]o0.north west) -- ([yshift=0.12cm]o2.north east) node[gruppe] {Patient};
    \draw[klammer] ([yshift=0.12cm]o3.north west) -- ([yshift=0.12cm]o4.north east) node[gruppe] {Probe};
    \draw[klammer] ([yshift=0.12cm]o5.north west) -- ([yshift=0.12cm]o8.north east) node[gruppe] {Messung};
  \end{tikzpicture}

  \vspace{0.45cm}
  \begin{columns}[t]
    \begin{column}{0.31\textwidth}
      \begin{kasten}
        \kastentitel{Patient}
        \small Verknüpft die Matrix mit der klinischen Tabelle.
      \end{kasten}
    \end{column}
    \begin{column}{0.31\textwidth}
      \begin{kasten}
        \kastentitel{Probentyp}
        \small Unterscheidet Primärtumor, Metastase und Normalgewebe.
      \end{kasten}
    \end{column}
    \begin{column}{0.31\textwidth}
      \begin{kasten}
        \kastentitel{Platte}
        \small Gemeinsam verarbeitet: Quelle von Batch-Effekten.
      \end{kasten}
    \end{column}
  \end{columns}
\end{frame}

% ------------------------------------------------------------------
\begin{frame}{Vom Gewebe zur Zahl: die Messung}
  \centering
  \begin{tikzpicture}[>={Stealth[length=5pt]},
      schritt/.style={font=\footnotesize\bfseries,text=blau,anchor=south},
      erkl/.style={font=\scriptsize,text=grau,align=center,anchor=north},
      pfeil/.style={->,grau,line width=0.9pt},
      read/.style={line width=2.2pt,line cap=round}]
    % 1 Gewebe
    \begin{scope}[shift={(0,0)}]
      \fill[petrolhell,draw=petrol!60,line width=0.8pt]
        plot[smooth cycle,tension=0.8] coordinates {(-1.0,0.2) (-0.6,0.85) (0.3,0.95) (1.0,0.45) (0.9,-0.5) (0.2,-0.95) (-0.8,-0.7)};
      \foreach \x/\y in {-0.55/0.35,-0.1/0.55,0.4/0.5,0.6/0.0,0.15/0.05,-0.35/-0.15,-0.6/-0.45,0.0/-0.5,0.45/-0.55,-0.75/0.0}
        {\draw[petrol!50,fill=white] (\x,\y) circle (0.19); \fill[petrol!70] (\x,\y) circle (0.06);}
      \node[schritt] at (0,1.25) {Gewebeprobe};
      \node[erkl] at (0,-1.25) {ein Stück Tumor,\\tiefgefroren};
    \end{scope}
    \draw[pfeil] (1.2,0) -- (1.75,0);
    % 2 RNA
    \begin{scope}[shift={(2.9,0)}]
      \foreach \y/\l in {0.7/1.6,0.35/1.1,0/1.8,-0.35/1.3,-0.7/1.5}
        \draw[blau,line width=1pt,domain=-0.9:-0.9+\l,samples=40,smooth] plot (\x,{\y+0.08*sin(\x*900)});
      \node[schritt] at (0,1.25) {RNA herauslösen};
      \node[erkl] at (0,-1.25) {Abschriften der gerade\\aktiven Gene};
    \end{scope}
    \draw[pfeil] (4.05,0) -- (4.6,0);
    % 3 Sequenzieren
    \begin{scope}[shift={(5.75,0)}]
      \node[font=\tiny\ttfamily,text=blau,align=left] at (0,0.05) {ACGTTAGCAT\\TTGACCGTAA\\GGCATTACGT\\CATGGTACCA\\ATTCGGATCG\\\dots};
      \node[schritt] at (0,1.25) {Sequenzieren};
      \node[erkl] at (0,-1.25) {Millionen kurzer\\Stücke (Reads) lesen};
    \end{scope}
    \draw[pfeil] (6.75,0) -- (7.3,0);
    % 4 Zuordnen
    \begin{scope}[shift={(8.85,0)}]
      \draw[grau,line width=0.8pt] (-1.45,-0.55) -- (1.45,-0.55);
      \foreach \x/\g in {-0.9/A,0.1/B,1.0/C}
        {\fill[blaurand] (\x-0.38,-0.7) rectangle (\x+0.38,-0.4);
         \node[font=\tiny\bfseries,text=white] at (\x,-0.55) {\g};}
      \foreach \x/\y in {-1.1/-0.25,-0.75/-0.25,-1.0/-0.05,-0.7/-0.05,-0.95/0.15,-0.8/0.35,-1.05/0.55,-0.7/0.55}
        \draw[read,orange] (\x-0.17,\y) -- (\x+0.17,\y);
      \foreach \x/\y in {0.0/-0.25,0.2/-0.05}
        \draw[read,orange] (\x-0.17,\y) -- (\x+0.17,\y);
      \node[font=\tiny,text=grau,anchor=west] at (-1.5,-0.85) {Erbgut};
      \node[schritt] at (0,1.25) {Gene zuordnen};
      \node[erkl] at (0,-1.25) {jeden Read an seine\\Stelle im Erbgut legen};
    \end{scope}
    \draw[pfeil] (10.4,0) -- (10.95,0);
    % 5 Zählen
    \begin{scope}[shift={(11.95,0)}]
      \node[draw=blaurand,fill=blauhell,rounded corners=2pt,inner sep=4pt,font=\footnotesize] at (0,0) {%
        \begin{tabular}{@{}lr@{}}
          Gen A & \textbf{8}\\
          Gen B & \textbf{2}\\
          Gen C & \textbf{0}
        \end{tabular}};
      \node[schritt] at (0,1.25) {Zählen};
      \node[erkl] at (0,-1.25) {Reads pro Gen\\ergeben die Zahl};
    \end{scope}
  \end{tikzpicture}

  \vspace{0.6cm}
  \begin{kasten}
    \centering Je aktiver ein Gen, desto mehr Abschriften und Reads, desto \textbf{größer die Zahl}.
  \end{kasten}
\end{frame}

% ------------------------------------------------------------------
\begin{frame}{Das Problem: unterschiedliche Sequenziertiefe}
  \centering
  \begin{tikzpicture}[>={Stealth[length=5pt]},
      pfeil/.style={->,grau,line width=0.9pt},
      titel/.style={font=\normalsize\bfseries,text=blau},
      erkl/.style={font=\scriptsize,text=grau,align=center}]
    \foreach \px/\name/\tiefe/\fa/\fb/\wa/\wb in {0/Probe 1/{20 Millionen\\Reads}/0.8/0.2/800/200,
                                                    7.8/Probe 2/{40 Millionen\\Reads}/1.6/0.4/1.600/400} {
      \begin{scope}[shift={(\px,0)}]
        \node[titel] at (0,1.9) {\name};
        % Gewebe
        \begin{scope}[shift={(-2.9,0.3)},scale=0.6]
          \fill[petrolhell,draw=petrol!60,line width=0.8pt]
            plot[smooth cycle,tension=0.8] coordinates {(-1.0,0.2) (-0.6,0.85) (0.3,0.95) (1.0,0.45) (0.9,-0.5) (0.2,-0.95) (-0.8,-0.7)};
          \foreach \x/\y in {-0.55/0.35,-0.1/0.55,0.4/0.5,0.6/0.0,0.15/0.05,-0.35/-0.15,-0.6/-0.45,0.0/-0.5,0.45/-0.55,-0.75/0.0}
            {\draw[petrol!50,fill=white] (\x,\y) circle (0.19); \fill[petrol!70] (\x,\y) circle (0.06);}
        \end{scope}
        \node[erkl] at (-2.9,-0.45) {gleiches Gewebe};
        \draw[pfeil] (-2.1,0.3) -- (-0.75,0.3);
        \node[font=\scriptsize\bfseries,text=text,anchor=south,align=center] at (-1.42,0.35) {\tiefe};
        % Balken
        \draw[grau,line width=0.6pt] (-0.55,-0.5) -- (2.75,-0.5);
        \foreach \i/\g/\h/\w in {0/A/\fa/\wa,1/B/\fb/\wb,2/C/0/0} {
          \fill[blau] (0.1+1.05*\i-0.3,-0.5) rectangle (0.1+1.05*\i+0.3,-0.5+\h);
          \node[font=\footnotesize\bfseries,text=blau,anchor=south] at (0.1+1.05*\i,-0.5+\h) {\w};
          \node[font=\footnotesize,anchor=north] at (0.1+1.05*\i,-0.55) {Gen \g};
        }
      \end{scope}
    }
  \end{tikzpicture}

  \vspace{0.5cm}
  \begin{kasten}
    \centering Gleiche Biologie, aber \textbf{doppelte Zahlen}: Probe 2 sähe doppelt so aktiv aus.
  \end{kasten}
\end{frame}

% ------------------------------------------------------------------
\begin{frame}{Normalisierung: an einem Beispiel}
  \centering
  {\small Die beiden Proben von eben: Probe 2 wurde \textbf{doppelt so tief} sequenziert.}

  \vspace{0.1cm}
  \begin{tikzpicture}[>={Stealth[length=5pt]},
      tab/.style={draw=blaurand,fill=blauhell,rounded corners=2pt,inner sep=5pt,font=\footnotesize},
      ergebnis/.style={draw=petrol!70,fill=petrolhell,rounded corners=2pt,inner sep=5pt,font=\footnotesize},
      titel/.style={font=\footnotesize\bfseries,text=blau,anchor=south},
      erkl/.style={font=\scriptsize,text=grau,align=center,anchor=north},
      pfeil/.style={->,grau,line width=0.9pt}]
    \node[tab] (roh) at (0,0) {%
      \begin{tabular}{@{}lrr@{}}
                & Probe 1 & Probe 2\\\midrule
        Gen A   & 800     & 1.600\\
        Gen B   & 200     & 400\\
        Gen C   & 0       & 0\\
        \midrule
        \textbf{oberes Quartil} & \textbf{400} & \textbf{800}
      \end{tabular}};
    \node[titel] at (roh.north) {gezählte Reads};
    \node[erkl] at ([yshift=-0.12cm]roh.south) {Probe 2 hat überall doppelte Werte,\\nur weil mehr gelesen wurde};

    \node[draw=grau!60,rounded corners=2pt,fill=white,align=center,inner sep=6pt,font=\footnotesize] (op) at (5.0,0)
      {jeden Wert\\$\div$ \textbf{oberes Quartil}\\$\times$ \textbf{1.000}};
    \draw[pfeil] (roh.east) -- (op.west);

    \node[ergebnis] (norm) at (9.6,0) {%
      \begin{tabular}{@{}lrr@{}}
                & Probe 1 & Probe 2\\\midrule
        Gen A   & 2.000   & 2.000\\
        Gen B   & 500     & 500\\
        Gen C   & 0       & 0\\
        \midrule
        \textbf{oberes Quartil} & \textbf{1.000} & \textbf{1.000}
      \end{tabular}};
    \draw[pfeil] (op.east) -- (norm.west);
    \node[titel,text=petrol] at (norm.north) {normalisierte Werte};
    \node[erkl] at ([yshift=-0.12cm]norm.south) {beide Proben jetzt\\direkt vergleichbar};
  \end{tikzpicture}

  \vspace{-0.25cm}
  \begin{columns}[t]
    \begin{column}{0.48\textwidth}
      \begin{kasten}
        \kastentitel{Oberes Quartil}
        \small Der Wert, unter dem 75 Prozent der Gene liegen.
      \end{kasten}
    \end{column}
    \begin{column}{0.48\textwidth}
      \begin{kasten}
        \kastentitel{Warum nicht die Summe?}
        \small Wenige extrem aktive Gene würden sie bestimmen.
      \end{kasten}
    \end{column}
  \end{columns}
\end{frame}

% ------------------------------------------------------------------
\begin{frame}{Einen Wert lesen}
  \centering
  \begin{tikzpicture}[>={Stealth[length=5pt]},
      zelle/.style={draw=blaurand,fill=blauhell,minimum width=1.5cm,minimum height=0.85cm,align=center,
                    font=\normalsize\bfseries,text=blau},
      hervor/.style={draw=petrol,fill=petrolhell,text=petrol,line width=1.2pt},
      kopf/.style={minimum width=1.5cm,font=\footnotesize\bfseries,anchor=south},
      zeile/.style={anchor=east,font=\footnotesize}]
    \def\w{1.5}
    \node[zeile] at (-0.8,0) {TCGA-05-4402-\dots};
    \foreach \i/\gen/\wert/\stil in {0/GAPDH/85.894/,1/XIST/3.458/,2/RPS4Y1/{1,2}/,3/SFTPC/1.610/,4/NKX2-1/2.779/hervor,5/ALB/{25,6}/} {
      \node[zelle,\stil] (z\i) at (\w*\i,0) {\wert};
      \node[kopf] at (z\i.north) {\gen};
    }
    \node[zelle,minimum width=0.9cm,draw=blaurand!60,dashed,fill=white,text=grau,font=\footnotesize] at (8.7,0) {\dots};
    % Rangskala
    \def\L{9.6}
    \begin{scope}[shift={(-0.75,-2.0)}]
      \fill[blauhell] (0,-0.12) rectangle (\L,0.12);
      \draw[blaurand] (0,-0.12) rectangle (\L,0.12);
      \foreach \p in {0,25,50,100}
        \node[font=\tiny,text=grau,anchor=north] at (\L*\p/100,-0.15) {\p\,\%};
      \foreach \p/\gen in {6.2/RPS4Y1,22.4/ALB,84.9/SFTPC,100/GAPDH}
        {\fill[blau] (\L*\p/100,0) circle (0.08);
         \node[font=\scriptsize,text=blau,above=2pt] at (\L*\p/100,0.1) {\gen};}
      \draw[orange,line width=1.2pt] (\L*0.75,-0.25) -- (\L*0.75,0.25);
      \node[font=\scriptsize\bfseries,text=orange,anchor=south] at (\L*0.75,0.25) {1.000};
      \node[font=\tiny,text=orange,anchor=north] at (\L*0.75,-0.15) {75\,\% = oberes Quartil};
      \fill[petrol] (\L*92.2/100,0) circle (0.11);
      \node[font=\scriptsize\bfseries,text=petrol,below=14pt] (nk) at (\L*92.2/100,0) {NKX2-1: 92\,\%};
      \node[font=\scriptsize,text=grau,anchor=east,align=right] at (-0.2,0) {Anteil der Gene über 0\\[-2pt] mit kleinerem Wert};
    \end{scope}
    \draw[->,petrol,line width=0.9pt] (z4.south) .. controls +(0,-0.6) and +(0,0.7) .. (-0.75+9.6*0.922,-1.85);
  \end{tikzpicture}

  \vspace{0.25cm}
  \begin{columns}[t]
    \begin{column}{0.48\textwidth}
      \begin{kasten}
        \kastentitel{Relativ zur Probe}
        \small 2.779 $\div$ 1.000 $\approx$ 2,8: NKX2-1 ist 2,8-mal so aktiv wie ein Gen am oberen Quartil, aktiver als 92 Prozent der Gene mit einem Wert über 0.
      \end{kasten}
    \end{column}
    \begin{column}{0.48\textwidth}
      \begin{kasten}
        \kastentitel{Biologisch}
        \small NKX2-1 steuert die Entwicklung der Lunge: ein hoher Wert passt zum Adenokarzinom der Lunge.
      \end{kasten}
    \end{column}
  \end{columns}
\end{frame}

% ------------------------------------------------------------------

% ==================================================================
\abschnittsfolie{2.\enspace Die Aufbereitung}

% ------------------------------------------------------------------ Korrekturen der Autoren
\begin{frame}{Was die Autoren schon korrigiert haben}
  \centering
  \vspace*{0.25cm}
  \begin{tikzpicture}[>={Stealth[length=5pt]},
      schritt/.style={draw=blaurand,fill=blauhell,rounded corners=3pt,minimum width=4.0cm,inner xsep=0.1cm,inner ysep=4pt,
                      minimum height=1.55cm,align=center},
      nr/.style={circle,fill=blau,text=white,font=\scriptsize\bfseries,inner sep=1.5pt},
      pfeil/.style={->,grau,line width=0.9pt}]
    \node[schritt] (s1) at (-0.15,0) {{\footnotesize\bfseries Normalisierung}\\[2pt]
      {\scriptsize alle Proben auf}\\{\scriptsize oberes Quartil = 1.000}};
    \node[schritt] (s2) at (4.35,0) {{\footnotesize\bfseries Batch-Korrektur (EB++)}\\[2pt]
      {\scriptsize Gerät und Zentrum}\\{\scriptsize pro Gen herausgerechnet}};
    \node[schritt] (s3) at (8.85,0) {{\footnotesize\bfseries Qualitätsprüfung}\\[2pt]
      {\scriptsize „nicht verwenden“ markiert:}\\{\scriptsize 184 Proben, nicht entfernt}};
    \foreach \s/\n in {s1/1,s2/2,s3/3} \node[nr] at (\s.north west) {\n};
    \draw[pfeil] (s1) -- (s2);
    \draw[pfeil] (s2) -- (s3);

    % Welche Proben hat EB++ verändert? (Breite proportional zur Probenzahl, 11.069 = 13 cm)
    \begin{scope}[shift={(-2.15,-2.45)}]
      \node[font=\footnotesize\bfseries,anchor=south west] at (0,0.42) {Welche Proben hat die Batch-Korrektur verändert?};
      \fill[blauhell] (0,-0.3) rectangle (10.775,0.3);
      \fill[orange!75] (10.775,-0.3) rectangle (11.477,0.3);
      \fill[orange] (11.477,-0.3) rectangle (12.763,0.3);
      \fill[orange!50] (12.763,-0.3) rectangle (12.802,0.3);
      \fill[grau!45] (12.802,-0.3) rectangle (13.0,0.3);
      \draw[blaurand] (0,-0.3) rectangle (13.0,0.3);
      \node[font=\footnotesize,text=blau] at (5.39,0) {\textbf{9.174} Illumina HiSeq, Hauptzentrum: \textbf{Referenz, unverändert}};
      \foreach \x/\farbe/\text in {0/orange!75/{\textbf{598} älteres Gerät\\{\color{grau}korrigiert}},
                                    4.5/orange/{\textbf{1.095} andere Zentren\\{\color{grau}korrigiert}},
                                    9.0/grau!45/{\textbf{169} Glioblastom (Broad Institute)\\{\color{grau}nicht korrigiert}}} {
        \fill[\farbe] (\x,-0.72) rectangle (\x+0.25,-0.47);
        \node[font=\scriptsize,anchor=north west,align=left,inner sep=0pt] at (\x+0.35,-0.45) {\text};
      }
    \end{scope}
    % zwei Kästen: gleiche Breite und Höhe, Abstand 0,5 cm wie zwischen den Schritten
    \node[anchor=north,inner sep=0pt] at (0.975,-4.3) {\begin{kasten}[width=6.25cm,equal height group=autoren]
        \kastentitel{Nicht korrigiert}
        \small Platten und Kliniken: Batch-Effekte können übrig sein.
      \end{kasten}};
    \node[anchor=north,inner sep=0pt] at (7.725,-4.3) {\begin{kasten}[width=6.25cm,equal height group=autoren]
        \kastentitel{Nebenwirkungen}
        \small Fehlende und leicht negative Werte in den korrigierten Proben.
      \end{kasten}};
  \end{tikzpicture}
\end{frame}

% ------------------------------------------------------------------ Probentypen
\begin{frame}{Welche Proben gibt es?}
  \bild[0.9\textwidth]{probentypen.pdf}
\end{frame}

\begin{frame}{Welche Proben behalten wir?}
  \bild[0.9\textwidth]{probentypen.pdf}
  \vspace{0.1cm}
  \begin{entscheidung}[fonttitle=\scriptsize\bfseries, toptitle=0.5pt, bottomtitle=0.5pt, top=2pt, bottom=2pt]
    {\normalsize\bfseries\color{petrol}Nur Primärtumoren}\quad{\small\color{grau}11.069 $\rightarrow$ 9.879 Proben}
  \end{entscheidung}
\end{frame}

% ------------------------------------------------------------------ Sequenziergerät

% ------------------------------------------------------------------ Fehlende Werte
\begin{frame}{Fehlende Werte}
  \centering
  \vspace*{0.3cm}
  \small
  \begin{tabular}{@{}l r r@{}}
    \toprule
    \textbf{Tumortyp} & \textbf{betroffene Proben} & \textbf{Gene mit fehlenden Werten} \\
    \midrule
    Leukämie (LAML)          & 100\,\% & 3.766 \\
    Magen (STAD)             & 100\,\% & 3.766 \\
    Dickdarm (COAD)          & 37\,\%  & 3.024 \\
    Mastdarm (READ)          & 42\,\%  & 3.024 \\
    Gebärmutterkörper (UCEC) & 67\,\%  & 3.024 \\
    Speiseröhre (ESCA)       & 100\,\% & 1.455 \\
    Eierstock (OV)           & 100\,\% & 1.455 \\
    Prostata (PRAD)          & 7\,\%   & 1.443 \\
    \midrule
    \textbf{zusammen}        &         & \textbf{4.196} \\
    \bottomrule
  \end{tabular}
\end{frame}

\begin{frame}{Was tun mit den fehlenden Werten?}
  \begin{entscheidung}
    \gross{Gene mit fehlenden Werten entfernen}\quad{\color{grau}20.531 $\rightarrow$ 16.335 Gene}
  \end{entscheidung}
  \vspace{0.2cm}
  \begin{columns}[t]
    \begin{column}{0.5\textwidth}
      \begin{itemize}
        \item Proben entfernen? Magen, Speiseröhre, Eierstock und Leukämie fielen ganz weg.
        \item Auffüllen? Ganze Tumortypen ohne einen einzigen Wert: Jede Auffüllung wäre erfunden.
      \end{itemize}
    \end{column}
    \begin{column}{0.46\textwidth}
      \begin{kasten}
        \kastentitel{Geprüft}
        \small Mit 0 aufgefüllt bleiben 1.599 wichtige Gene erhalten (zum Beispiel Keratine).
        Nachbarn aus demselben Tumortyp: \textbf{85,0\,\%} statt \textbf{84,3\,\%}. Kaum ein Unterschied.
      \end{kasten}
    \end{column}
  \end{columns}
\end{frame}

% ------------------------------------------------------------------ Log-Transformation
\begin{frame}{Die Werte vor der Transformation}
  \bild{werte_roh.pdf}
  \vspace{0.15cm}
  \begin{kasten}
    Stark \textbf{rechtsschief}. Wenige sehr hoch exprimierte Gene tragen fast die ganze Varianz:
    Jede Abstands- und Varianzrechnung würde nur sie sehen.
  \end{kasten}
\end{frame}

\begin{frame}{Transformation}
  \begin{entscheidung}
    \gross{$x \mapsto \log_2(x + 1)$}\quad{\color{grau}negative Werte vorher auf 0}
  \end{entscheidung}
  \vspace{0.3cm}
  \begin{itemize}
    \item Ein Unterschied von 1 heißt \textbf{Verdopplung}, egal auf welchem Niveau.
    \item $+1$: damit $0$ definiert bleibt ($\log_2 1 = 0$).
    \item Negative Werte (1,25 Millionen, alle zwischen $-1$ und $0$): Nebenwirkung der Batch-Korrektur, bedeuten „nicht exprimiert“.
  \end{itemize}
\end{frame}

\begin{frame}{Die Werte nach der Transformation}
  \bild{werte_log.pdf}
  \vspace{0.15cm}
  \begin{kasten}
    Annähernd symmetrisch (plus eine Spitze bei 0: nicht exprimiert). Die Varianz verteilt sich auf viele Gene:
    Die 50 stärksten tragen nur noch \textbf{3\,\%} statt \textbf{81\,\%}.
  \end{kasten}
\end{frame}

% ------------------------------------------------------------------ Genauswahl
\begin{frame}{Wie viel streuen die Gene?}
  \bild[0.78\textwidth]{varianz_gene.pdf}
  \vspace{0.1cm}
  \begin{kasten}
    Wenige Gene streuen stark, die meisten kaum. Die 5.000 Gene mit der größten Varianz tragen \textbf{72\,\%} der Gesamtvarianz.
  \end{kasten}
\end{frame}

\begin{frame}{Welche Gene behalten wir?}
  \begin{entscheidung}
    \gross{Die 5.000 Gene mit der größten Varianz}\quad{\color{grau}16.331 $\rightarrow$ 5.000 Gene}
  \end{entscheidung}
  \vspace{0.2cm}
  \begin{itemize}
    \item Varianz auf der \textbf{logarithmierten} Skala, sonst gewinnen immer die hoch exprimierten Gene.
    \item Vorher: Gene entfernen, die in keinem Tumortyp bei mindestens der Hälfte der Proben exprimiert sind (trifft nach dem vorigen Schritt kein Gen mehr).
    \item Pro Gen \textbf{zentrieren}, nicht standardisieren: stark streuende Gewebegene sollen die Achsen prägen.
  \end{itemize}
\end{frame}

% ------------------------------------------------------------------ Geschlechtsgene
\begin{frame}{Welche Gene hängen am Geschlecht?}
  \bild[0.85\textwidth]{geschlecht.pdf}
  \vspace{0.1cm}
  \begin{kasten}
    Nur \textbf{4 Gene} hängen stark am Geschlecht: \gen{XIST} und \gen{TSIX} (nur bei Frauen aktiv) und zwei Gene des Y-Chromosoms.
    Die übrigen Y-Gene sind schon mit den fehlenden Werten weggefallen. Danach eine klare Lücke.
  \end{kasten}
\end{frame}

\begin{frame}{Geschlechtsgene}
  \begin{entscheidung}
    \gross{4 Geschlechtsgene entfernen}\quad{\color{grau}aus den Daten bestimmt: $r^2 > 0{,}4$}
  \end{entscheidung}
  \vspace{0.3cm}
  \begin{itemize}
    \item Gefragt ist der Tumortyp, nicht das Geschlecht der Patienten.
    \item $r^2$ wird \textbf{innerhalb} der Tumortypen gemessen: Prostata-Gene zählen nicht als Geschlechtsgene.
    \item Danach erklärt das Geschlecht auf keiner der ersten 20 Hauptkomponenten mehr als 0,1\,\%.
  \end{itemize}
\end{frame}

% ------------------------------------------------------------------ Batch-Effekte
\begin{frame}{Was prägt die Hauptkomponenten?}
  \bild{stoergroessen.pdf}
  \vspace{0.1cm}
  \begin{kasten}
    Der \textbf{Tumortyp} erklärt den Großteil fast jeder Achse. Die Platte (Batch) erklärt etwas mehr als der Zufall,
    am meisten auf den Hauptkomponenten 8, 11 und 12.
  \end{kasten}
\end{frame}

\begin{frame}{Platteneffekt innerhalb des Brustkrebses}
  \bild[0.88\textwidth]{platte_brustkrebs.pdf}
  \vspace{0.1cm}
  \begin{kasten}
    Eigene Hauptkomponentenanalyse nur für Brustkrebs: Achse 3 wird zu \textbf{51\,\%} durch die Platte erklärt (Zufall: 4\,\%).
    Die Platten verschieben sich schrittweise, vermutlich mit dem Zeitpunkt der Verarbeitung.
  \end{kasten}
\end{frame}

\begin{frame}{Batch-Effekte}
  \begin{entscheidung}
    \gross{Keine zusätzliche Korrektur, aber kontrollieren}
  \end{entscheidung}
  \vspace{0.3cm}
  \begin{itemize}
    \item Platten und Zentren sind mit dem Tumortyp vermengt: 55\,\% der Platten enthalten nur einen Tumortyp.
          Eine Korrektur nach Platte würde echte Unterschiede mit entfernen.
    \item Auffällige Gruppen später zuerst nach Platte einfärben.
  \end{itemize}
\end{frame}

% ------------------------------------------------------------------ Ergebnis
\begin{frame}{Ergebnis der Aufbereitung}
  \bild{protokoll.pdf}
  \vspace{0.15cm}
  \begin{kasten}
    \textbf{9.720 Primärtumoren $\times$ 5.000 Gene}, logarithmiert und zentriert. Geprüft und verworfen:
    Ausgleich der Gruppengrößen und weitere Qualitätsfilter. Beides ändert die Nachbarschaften nicht.
  \end{kasten}
\end{frame}

% ==================================================================
\abschnittsfolie{3.\enspace Hauptkomponenten\-analyse}

\begin{frame}{Idee}
  \begin{columns}[c]
    \begin{column}{0.42\textwidth}
      \centering
      \begin{tikzpicture}[>={Stealth[length=5pt]}]
        \pgfmathsetseed{4}
        \begin{scope}[rotate=28]
          \foreach \i in {1,...,70} {
            \pgfmathsetmacro{\a}{2.0*(rnd+rnd+rnd-1.5)}
            \pgfmathsetmacro{\b}{0.7*(rnd+rnd+rnd-1.5)}
            \fill[blaurand,opacity=0.8] (\a,\b) circle (1.3pt);
          }
          \draw[->,petrol,line width=1.3pt] (0,0) -- (2.4,0) node[right,font=\footnotesize] {1};
          \draw[->,orange!85!black,line width=1.3pt] (0,0) -- (0,1.1) node[above,font=\footnotesize] {2};
        \end{scope}
        \node[font=\scriptsize,text=grau] at (0,-2.0) {Punkte: Proben; Pfeile: Hauptkomponenten};
      \end{tikzpicture}
    \end{column}
    \begin{column}{0.54\textwidth}
      \begin{itemize}
        \item Neue Achsen in Richtung der \textbf{größten Streuung}, senkrecht zueinander.
        \item Jede Achse ist eine gewichtete Summe aller Gene (\textbf{Ladungen}).
        \item Berechnung: abgeschnittene Singulärwertzerlegung $X \approx U_k \Sigma_k V_k^\top$, nur die ersten 100 Achsen.
      \end{itemize}
    \end{column}
  \end{columns}
\end{frame}

\begin{frame}{Wie viel Varianz erklären die Hauptkomponenten?}
  \bild{pca_varianz.pdf}
  \vspace{0.1cm}
  \begin{kasten}
    Statt 5.000 Genen genügen wenige Dutzend Achsen: 10 erklären 53\,\%, 50 rund 72\,\%.
    Ab etwa 15 bis 20 flacht die Kurve ab, aber ohne klaren Knick.
  \end{kasten}
\end{frame}

\begin{frame}{Hauptkomponenten 1 und 2}
  \bild[0.72\textwidth]{pca_12.pdf}
  \vspace{-0.05cm}
  \begin{kasten}
    \textbf{Achse 1:} Gehirn und Nebenniere gegen Epithel. \textbf{Achse 2:} Leber und Leukämie heraus.
    Die Karzinome überlappen links oben in einer großen Wolke.
  \end{kasten}
\end{frame}

\begin{frame}{Hauptkomponenten 3 und 4}
  \bild[0.72\textwidth]{pca_34.pdf}
  \vspace{-0.05cm}
  \begin{kasten}
    \textbf{Achse 3:} Nieren links, Leukämie, Lymphom und Melanome rechts. \textbf{Achse 4:} Immunzellen und Leukämie oben, Prostata unten.
  \end{kasten}
\end{frame}

\begin{frame}{Was bedeuten die Achsen?}
  \bild[0.92\textwidth]{pca_ladungen.pdf}
  \quelle{Die 5 Gene mit den größten positiven und negativen Ladungen. Das Vorzeichen einer Achse ist willkürlich.}
  \vspace{0.05cm}
  \begin{kasten}
    Die Achsen sind \textbf{Gewebeprogramme}: Nervengewebe (\gen{GFAP}, \gen{PLP1}), Leber (\gen{ALDOB}, \gen{PKLR}),
    Plattenepithel (\gen{KRT17}), Drüsen (\gen{AGR2}), Immunzellen (\gen{S100A8}, \gen{GZMK}).
  \end{kasten}
\end{frame}

\begin{frame}{Zwischenfazit Hauptkomponentenanalyse}
  \begin{columns}[t]
    \begin{column}{0.48\textwidth}
      \begin{kasten}
        \kastentitel{Ergebnis}
        \small
        \begin{itemize}
          \item Die Expression ist vom \textbf{Herkunftsgewebe} geprägt.
          \item Klar abgesetzt: Gehirn, Leber, Leukämie, Niere, Melanome, Prostata.
          \item Die meisten Karzinome überlappen; ihre Unterschiede verteilen sich auf viele Achsen.
        \end{itemize}
      \end{kasten}
    \end{column}
    \begin{column}{0.48\textwidth}
      \begin{entscheidung}
        \small Die ersten \textbf{50 Hauptkomponenten} (72\,\% der Varianz) sind die Eingabe für UMAP und das Clustering.
      \end{entscheidung}
    \end{column}
  \end{columns}
\end{frame}

% ==================================================================
\abschnittsfolie{4.\enspace UMAP}

\begin{frame}{Idee}
  \centering
  \begin{tikzpicture}[>={Stealth[length=4pt]},
                      p/.style={circle,inner sep=1.6pt},
                      titel/.style={font=\footnotesize\bfseries,text=text},
                      erkl/.style={font=\scriptsize,text=grau,text width=3.6cm,align=center}]
    % 1: Nachbarn
    \begin{scope}
      \foreach \n/\x/\y/\f in {a/0/0/blau,b/0.5/0.35/blau,c/0.35/-0.4/blau,d/0.9/-0.1/blau,
                               e/2.0/0.6/orange,f/2.5/0.2/orange,g/2.2/-0.3/orange,h/2.7/0.8/orange}
        \node[p,fill=\f] (\n) at (\x,\y) {};
      \foreach \u/\v in {a/b,a/c,b/d,c/d,a/d,e/f,f/g,e/h,f/h,e/g,d/g}
        \draw[grau!70] (\u) -- (\v);
      \node[titel] at (1.35,1.5) {1.\ Nachbarn suchen};
      \node[erkl] at (1.35,-1.25) {für jede Probe die 15 nächsten Nachbarn in den 50 Hauptkomponenten};
    \end{scope}
    \draw[->,grau,line width=0.8pt] (3.3,0.2) -- (4.1,0.2);
    % 2: Gewichte
    \begin{scope}[xshift=4.6cm]
      \foreach \n/\x/\y/\f in {a/0/0/blau,b/0.5/0.35/blau,c/0.35/-0.4/blau,d/0.9/-0.1/blau,
                               e/2.0/0.6/orange,f/2.5/0.2/orange,g/2.2/-0.3/orange,h/2.7/0.8/orange}
        \node[p,fill=\f] (\n) at (\x,\y) {};
      \foreach \u/\v/\w in {a/b/1.6,a/c/1.6,b/d/1.2,c/d/1.4,a/d/0.7,e/f/1.6,f/g/1.4,e/h/1.2,f/h/1.0,e/g/0.7,d/g/0.3}
        \draw[blau!70,line width=\w pt] (\u) -- (\v);
      \node[titel] at (1.35,1.5) {2.\ Kanten gewichten};
      \node[erkl] at (1.35,-1.25) {nahe Nachbarn stark, ferne schwach; relativ zur Dichte um jede Probe};
    \end{scope}
    \draw[->,grau,line width=0.8pt] (7.9,0.2) -- (8.7,0.2);
    % 3: Layout
    \begin{scope}[xshift=9.3cm]
      \foreach \x/\y in {0/0,0.35/0.3,0.3/-0.25,0.6/0.05} \node[p,fill=blau] at (\x,\y) {};
      \foreach \x/\y in {2.1/0.1,2.45/0.4,2.4/-0.2,2.75/0.15} \node[p,fill=orange] at (\x,\y) {};
      \draw[<->,grau] (0.8,0.05) -- (1.9,0.05);
      \node[titel] at (1.35,1.5) {3.\ in 2 Dimensionen legen};
      \node[erkl] at (1.35,-1.25) {Verbundene ziehen sich an, alle anderen stoßen sich ab};
    \end{scope}
  \end{tikzpicture}
  \vspace{0.5cm}
  \begin{kasten}
    UMAP erhält vor allem \textbf{lokale Nachbarschaften}. Abstände zwischen Gruppen und Gruppengrößen im Bild sind \textbf{nicht} deutbar.
  \end{kasten}
\end{frame}

\begin{frame}{Die Tumortypen im UMAP-Bild}
  \bild[0.8\textwidth]{umap.pdf}
  \quelle{Eingabe: 50 Hauptkomponenten; n\_neighbors = 15, min\_dist = 0,5. Beschriftung am Median jedes Tumortyps.}
  \vspace{0.05cm}
  \begin{kasten}
    \small\textbf{Abseits:} Tumortypen mit eigenem Gewebe (Leukämie, Prostata, Leber, Niere, Gehirn). \textbf{Mitte:} die Karzinome,
    aber in Gruppen (Plattenepithel, Verdauungstrakt, Gynäkologie). Die Karzinom-Wolke ist aufgelöst.
  \end{kasten}
\end{frame}

\begin{frame}{Wie klar sind die Tumortypen abgegrenzt?}
  \bild{eigener_anteil.pdf}
  \quelle{Für jede Probe die 15 nächsten Nachbarn in den 50 Hauptkomponenten, gemittelt pro Tumortyp.}
  \vspace{0.05cm}
  \begin{kasten}
    Die meisten Tumortypen sind \textbf{klar abgegrenzt}. Am wenigsten eigenständig: Mastdarm, Speiseröhre, Karzinosarkom, Gallengang.
  \end{kasten}
\end{frame}

\begin{frame}{Welche Tumortypen ähneln sich?}
  \begin{columns}[c]
    \begin{column}{0.56\textwidth}
      \bild{paare.pdf}
    \end{column}
    \begin{column}{0.42\textwidth}
      \small
      Anteil der Nachbarn von A im Tumortyp B:
      \begin{itemize}
        \item \textbf{gleiches Organ:} Mast- und Dickdarm sind praktisch derselbe Tumor.
        \item \textbf{gleiches Gewebe:} Plattenepithelkarzinome verschiedener Organe.
        \item \textbf{gleiche Herkunft:} Gallengang und Leber; Gebärmutter und Eierstock.
      \end{itemize}
    \end{column}
  \end{columns}
\end{frame}

\begin{frame}{Tumortypen, die in zwei Teile zerfallen}
  \bild[0.8\textwidth]{untertypen.pdf}
  \vspace{0.1cm}
  \begin{kasten}
    Die Teile sind \textbf{histologische Untertypen}: Plattenepithel bei Kopf-Hals und Lunge, Adenokarzinom beim Magen;
    Brust: Östrogenrezeptor hoch oder niedrig.
  \end{kasten}
\end{frame}

\begin{frame}{Wie stabil ist das Bild? Parameter}
  \bild[0.72\textwidth]{umap_parameter.pdf}
  \vspace{-0.05cm}
  \begin{kasten}
    Die Bilder sehen sehr verschieden aus, aber \textbf{dieselben Tumortypen liegen zusammen}.
    Anteil der Nachbarn aus demselben Tumortyp: 82 bis 83\,\% in allen sechs Bildern.
  \end{kasten}
\end{frame}

\begin{frame}{Wie stabil ist das Bild? Zufallsstartwert}
  \bild{umap_startwert.pdf}
  \vspace{0.1cm}
  \begin{kasten}
    Gleiche Daten, anderer Startwert: Prostata, Leber und Brustkrebs liegen jedes Mal woanders.
    Die \textbf{Lage der Gruppen} ist zufällig, ihre \textbf{Zusammensetzung} nicht.
  \end{kasten}
\end{frame}

\begin{frame}{Was das UMAP-Bild leistet}
  \centering
  \vspace{0.2cm}
  \kennzahl{84,3\,\%}{Nachbarn aus demselben Tumortyp\\in 50 Hauptkomponenten}\hfill
  \kennzahl{82,2\,\%}{Nachbarn aus demselben Tumortyp\\im UMAP-Bild}\hfill
  \kennzahl{34,1\,\%}{der Nachbarn einzelner Proben\\im Bild erhalten}
  \vspace{0.6cm}
  \raggedright
  \begin{kasten}
    Gute \textbf{Übersicht} über die Tumortypen, für einzelne Proben ungenau.
    Gemessen wird deshalb in den 50 Hauptkomponenten, nicht im Bild.
  \end{kasten}
\end{frame}

% ==================================================================
\begin{frame}{Fazit und nächste Schritte}
  \begin{columns}[t]
    \begin{column}{0.48\textwidth}
      \begin{kasten}
        \kastentitel{Bisher}
        \small
        \begin{itemize}
          \item Die meisten Tumortypen sind klar abgegrenzt.
          \item Ähnlichkeit folgt dem \textbf{Gewebe}: gleiches Organ, gleicher Gewebetyp, gleiche Herkunft.
          \item Bei Speiseröhre und Gebärmutterhals ist der Gewebetyp wichtiger als das Organ.
        \end{itemize}
      \end{kasten}
    \end{column}
    \begin{column}{0.48\textwidth}
      \begin{kasten}
        \kastentitel{Als Nächstes}
        \small
        \begin{itemize}
          \item Spektrales Clustering auf dem Nachbarschaftsgraphen und ein zweites Verfahren.
          \item Stimmen die Cluster mit den Tumortypen überein? Finden sie die Untertypen?
          \item Robustheit: Zahl der Gene, Standardisieren, Abstandsmaß.
        \end{itemize}
      \end{kasten}
    \end{column}
  \end{columns}
\end{frame}

\end{document}
